\chapter{Теоретична основа}
\label{ch:osnova}

\todo{Фоновите знания, нужни за разбиране на работата. По стила на ФМИ
свързаните разработки могат да се вплетат тук с цитати.}

\section{Барабанна динамика и възприятие}
\label{sec:dinamika}
\todo{Какво е velocity, как влияе на тембъра и усещането за „грув''.}

\section{MIDI и представяне на барабанни партии}
\label{sec:midi}
\todo{MIDI събития, барабанен маппинг, квантуване, timing.}

\section{Признаци от структура и време}
\label{sec:priznaci}
\todo{Метрична позиция, такт/доли, плътност, съседни ноти, инструмент и т.н.}

\section{Наборът от данни E-GMD}
\label{sec:egmd}
\todo{Описание на E-GMD: обхват, брой изпълнения, особености (виж бележките
от EDA и артефакта с пренасочването на гласове между китове).}
